\chapter*{Foreword}
\addcontentsline{toc}{chapter}{Foreword}
\markboth{Foreword}{Foreword}

The large \text{history of Soviet Russia} which has occupied me for the past thirty years, and has just been completed in four instalments, \textit{The Bolshevik Revolution, 1917-1923}, \textit{The Interregnum, 1923-1924}, \textit{Socialism in One Country, 1924-1926}, and \textit{Foundations of a Planned Economy, 1926-1929}, was based on much detailed research and designed for specialists. It occurred to me that some purpose might be served by distilling this research into a short book of a quite different kind, without the scholarly refinements of source references or footnotes, designed for the general reader and for the student seeking a first introduction to the subject. The result is the present short history. The difference in scale and purpose means that this is substantially a new composition. Scarcely a sentence from the original work reappears unchanged in the new. 

\textit{The Russian Revolution: from Lenin to Stalin, 1917-1929} covers the same period as the large history. This is a period for which (in contrast with the later years) ample contemporary Soviet sources are available. It is also a period which contained in embryo much of the subsequent course of Soviet history; an understanding of what happened then is needed to explain what happened afterwards. % 注：此处“afterwards”疑为“afterwards”的 OCR 错误。
To describe the nineteen-twenties in terms of a transition from the Russian revolution of Lenin to the Russian revolution of Stalin would, no doubt, be an over-simplification. But it would personify an important historical process, the conclusion of which still lies in the unforeseeable future. 

The many friends and colleagues whose names are cited in the prefaces to successive volumes of the large history also claim acknowledgment here as indirect contributors to the present work. To Professor R. W. Davies, who collaborated with me in the first volume of \textit{Foundations of a Planned Economy, 1926-1929}, I am specially indebted for expert criticism of the chapters on industrialization and planning; and I have read with profit Professor Alec Nove's concise \textit{Economic History of the USSR}. My warm thanks are once more due to Tamara Deutscher for her unfailing help in the preparation of this volume. 

November 7 1977 % 注：此处“igyy”疑为年份“1977”。
\\
E. H. Carr 


\chapter*{List of Abbreviations}
\addcontentsline{toc}{chapter}{List of Abbreviations}
\markboth{List of Abbreviations}{List of Abbreviations}

% 使用 longtable 制作两列对比的表格，左侧宽度缩进，右侧自适应
\begin{longtable}{@{}p{3.5cm} p{10cm}@{}}
Arcos & All-Russian Cooperative Society \\
CCP & Chinese Communist Party \\
CER & Chinese Eastern Railway \\
Cheka & Extraordinary Commission \\
Comintern & Communist International \\
CPGB & Communist Party of Great Britain \\
Glavk(i) & Chief Committee(s) \\
Goelro & State Commission for the Electrification of Russia \\
Gosplan & State Planning Commission \\
GPU & State Political Administration \\
IFTU & International Federation of Trade Unions \\
IKKI & Executive Committee of the Communist International \\
Kadet & Constitutional-Democrat \\
Khozraschet & Commercial Accounting \\
Kolkhoz(y) & Collective Farm(s) \\
KPD & German Communist Party \\
MTS & Machine Tractor Station \\
Narkomfin & People's Commissariat of Finance \\
Narkomindel & People's Commissariat of Foreign Affairs \\
Narkomprod & People's Commissariat of Supply \\
Narkomtrud & People's Commissariat of Labour \\
Narkomzem & People's Commissariat of Agriculture \\
NEP & New Economic Policy \\
NMM & National Minority Movement \\
NUWM & National Unemployed Workers' Movement \\
OGPU & Unified State Politieal Administration \\ % 注：此处“Politieal”疑为“Political”。
PCF & French Communist Party \\
PCI & Italian Communist Party \\
Profintern & Red International of Trade Unions (RILU) \\
RSFSR & Russian Socialist Federal Soviet Republic \\
Sovkhoz(y) & Soviet Farm(s) \\
Sovnarkhoz(y) & Council(s) of National Economy \\
Sovnarkom & Council of People's Commissars \\
SPD & German Social-Democratic Party \\
SR & Social-Revolutionary \\
TsIK & Central Executive Committee \\
USPD & German Independent Social-Democratic Party \\
USSR & Union of Soviet Socialist Republics \\
VAPP & All-Russian Association of Proletarian Writers \\
Vesenkha & Supreme Council of National Economy \\
VTsIK & All-Russian Central Executive Committee \\
\end{longtable}